%%%%%%%%%%%%%%%%%%%%%%%%%%%%%%%%%%%%%%%%%%%%%%%%%%%%%
%			HLAVIČKA								%
%%%%%%%%%%%%%%%%%%%%%%%%%%%%%%%%%%%%%%%%%%%%%%%%%%%%%
\documentclass[openany,12pt]{memoir}

\usepackage[utf8]{inputenc} 
\usepackage[czech]{babel}
\usepackage[T1]{fontenc}
\usepackage[top=1.5cm, bottom=2cm, left=2cm, right=2cm]{geometry}  % --> NASTAVENÍ OKRAJŮ
\usepackage{fancyhdr}
\usepackage{graphicx}
%\usepackage{lmodern}
\usepackage{tgschola}  % --- Silnější font
\usepackage{xwatermark}
\usepackage{xcolor}
\usepackage{changepage}
\usepackage{newunicodechar}



%%%%%% Package na zpěvník
\usepackage[full]{leadsheets}%http://mirrors.nic.cz/tex-archive/macros/latex/contrib/leadsheets/leadsheets_en.pdf   --> dokumentace	
\definesongtitletemplate{empty}{} 
\setchords{
format = \bfseries,   %tučné akordy
minor = {mi},% 
input-notation = {german},%
output-notation = {german}%
}
\definesongtitletemplate{empty}{} 

\newlength{\drop}
\newwatermark[allpages,color=red!50,angle=0,scale=2, xpos=0,ypos=0]{\includegraphics[width=5cm]{obr/pozadi2.jpg}} %--> dvojka na pozadí


%%%% Vlastní příkazy
\newcounter{Slokočet}   %Automatické číslování slok
\newcommand{\mezera}{\vspace*{0.5cm}}   %Horizontální odsazení slok
\newcommand{\stred}{5.2cm}   %%% Na zarovnání slok doprostřed, pozn. automatičtější zarovnávání na střed nejde
\newcommand{\refren}{\mezera \noindent \textbf{R:} } %refrén
\newcommand{\sloka}{\mezera \noindent \addtocounter{Slokočet}{1} \arabic{Slokočet}. } 	%sloka, která se automaticky čísluje
\newcommand{\ssloka}{\mezera \noindent}  % vlastní číslo sloky
\newcommand{\carka}{,\:}
\newcommand{\predehra}{\mezera \hspace*{-1.85cm} \textbf{Předehra:} }
\newcommand{\m}[1]{\color{white}{#1}}  %Pro akordy
\newcommand{\ap}{'} %Pro apostrof
\newcommand{\elipsa}{\kern\fontdimen3\font} %Příkaz pro lepší zacházení s výpustkami (=...); je to vpodstatě jen mezera mezi tečkama výpustky

%%%%%%%%%%%%%%%%%%%%%%%%%%%%%%%%%%%%%%%%%%%%%%%%%%%%%%
%			NÁVOD									 %
%%%%%%%%%%%%%%%%%%%%%%%%%%%%%%%%%%%%%%%%%%%%%%%%%%%%%%
%1. Věci v hlavičce IGNOROVAT
%2. Píseň psát do prostoru mezi \begin{song} a \end{song}
%3. další řádek se značí dvěma odsazeníma (= dvakrát stisknout enter)
%4. \refren vždy na začátku refrenu a \sloka na začátek sloky (automaticky se čísluje)
%  \ = alt gr + q ; [/] = alt gr f/g ; {/} = alt gr + b/n; ^ = alt gr + 3 + mezera
%Cokoliv napíšete do ^{  } se bude brát jako akord
%když se toto bude dotýkat nějakého slova (nebude mezi tím a slovem mezera)
%tak se akordy zjeví nad slovem, ale pište to před slovo
%Když se to nedotýká slova, tak akord lítá ve vzduchu a vytiskne se větší mezera
%První možnost je asi preferovanější
%5. Akordy stačí psát jen do první sloky, když se nezmění -- kytaristi to zvládnou



%\usepackage{subfiles}
%</++>

\begin{document}

\pagestyle{empty}
%\addtocontents{toc}{\protect\thispagestyle{empty}}
%\tableofcontents \thispagestyle{empty}\newpage
\newgeometry{top=1.5cm, bottom = 0cm, left = 2cm, right = 2cm}

\newpage

%Pro kompilaci je důležitý, 	aby ten input a \end{document} zůstali na posledních
%dvou řádcích!!!!
<<<<<<< HEAD
\begin{song}{title=\predtitle\centering Kaliforňané \\\large Greenhorns  \vspace*{-0.3cm}}  %% sem se napíše jméno songu a autor
\begin{centerjustified}
\nejnejvetsi

\sloka
	Ať ^{C\z}všichni ^{Ami}Kaliforňané se ^{G}za nás ^{\z C}pomodlí,

	kdo táhneš s mulou ^{F\z}přes pláně, ^{\z C}rozluč se s ^{\z F}pohodlím,

   	^{C\z }musíš si taky ^{F\z}nakoupit ^{C}solený ^{\z F}vepřový,

   	a ^{C}kdo má štěstí, ^{Ami\z}po~cestě i ^{G\z}něco ^{\z C}uloví.

\sloka
	Po boku měj dva kolťáky a vostrej Bowie nůž.

   	Jestli tu cestu překonáš, potom seš celej muž.

   	Zelenáči, až uvidíš kol sebe poušť a prach

   	a nikde živý stvoření, dostaneš asi strach.

\sloka
	Dva měsíce ti potrvá ta cesta přes pláně,

   	koukej, ať vždycky z čeho máš vařit svý snídaně,

   	protože k čertu dojedeš při bídným vaření,

   	i tvoje mula musí žrát, jináč ti zcepení.

 \sloka
	A taky proti rudochům drž stráže za noci,

   	po dlouhejch denních pochodech propuknou nemoci,

   	horečka, to je špatná věc, ta lidi poráží,

   	a kdo je marod, nemůže stát v noci na stráži.

 \sloka
	Břitvu nech doma, do tváře se radši nedívej,

   	až přijdeš, budeš zarostlej a taky šedivej,

   	až tě tvá milá uvidí, uslyšíš asi pláč,

   	svý mámě řekneš: ,,Podívej, to je tvůj zelenáč!{}``

\sloka
	U Plat je nouze vo dřevo, když tě to nezdolá,

   	omáčku si tam zamícháš paznehtem z buvola,

   	u Plat je nouze vo klacek, zem je tam jalová,

   	i já bych teď už radši byl u Leopoldova. %mysu Hornova 

\end{centerjustified}
\setcounter{Slokočet}{0}
\end{song}

=======
\begin{song}{title=\centering One \\\normalsize Johnny Cash  \vspace*{-0.3cm}}  %% sem se napíše jméno songu a autor

\renewcommand{\baselinestretch}{0.4}

\moveright 0.0cm \vbox{      %Varianta č. 1  ---> Jeden sloupec zarovnaný na střed	

\begin{minipage}[t]{0.55\textwidth}\setlength{\parindent}{0.45cm} %Varianta č. 2 --> Dva sloupce


%\small

\sloka 
 	^{Ami}Is it getting ^{D{\color{white}\_\_\_}}better

	^{F\,}Or do you feel the ^{G{\color{white}\_\_\_}}same?

	^{Ami}Will it make it ^{D{\color{white}\_\_\_}}easier on you now

	^{F\,\,}If you've got someone to ^{{\color{white}\_\_\_}G}blame?

   
\refren
	You said ^{C}one love ^{Ami}

	One life ^{F}

	When it's one need ^{C}
         
	In the night.
  
	One love, we get to ^{Ami\,\,}share it.

	It ^{F{\color{white}\_\_}}leaves you, baby, if you don't ^{C{\color{white}\_\_}}care for it.   


\sloka
	Did I disappoint you

	Or leave a bad taste in your mouth?

	You act like you never had love

	And you want me to go without.


\refren
	Well it's too late 

	Tonight 

	To drag the past out 

	Into the light.

	We're one but we're not the same.

	We get to carry each other,

	Carry each other.
 
	One.

\sloka    
	Have you come here for forgiveness?

	Have you come to raise the dead?

	Have you come here to play Jesus

	To the lepers in your head?


\end{minipage}\begin{minipage}[t]{0.55\textwidth}\setlength{\parindent}{0.45cm}\vspace*{0.55cm}  % V případě varianty č.2 jde odsud text do pravé části


\refren       
	Did I ask too much? 

	More than a lot?

	You gave me nothing, now
	   
	It's all I got.

	We're one but we're not the same.

	Well, we hurt each other and we do it again. 
 
\refren
	You said ^{C{\color{white}\_}}love is a temple.

	^{Ami\,\,}Love the higher law.

	Love is a temple.	

	Love the higher law.

	^{C{\color{white}\_}}You ask me to enter

	^{G\,\,}But then you make me crawl

	I can't be holding on ^{F}
                  
	To what you've got

	When all you've got is ^{C}hurt.
  
\refren
	One love, ^{Ami}

	One blood, ^{F}

	One life, ^{C}
 
	You've got to do what you should.

	One life with ^{Ami\,\,}each other. ^{F}
 
 	Sisters, ^{C}
 
	Brothers.
   
	One life, but we're ^{Ami}not the same.

	We get to ^{F}carry each other,

  	^{C{\color{white}\_\_}}Carry each other,

  	/: ^{C\,\,Ami\,\,}One ^{F\,\,C\,\,}One. :/


\normalsize


\end{minipage}
}
\setcounter{Slokočet}{0}
\end{song}

>>>>>>> f53c6ab8e34a4a131bef5f220cd3dd45ab0906ee
\end{document}
